\documentclass[10pt,a4paper]{amsart}
\usepackage[margin=19mm]{geometry}
\usepackage{amsmath,amssymb,mathtools}
\usepackage[scr=rsfs,cal=euler]{mathalfa}
\usepackage{xfrac}
\usepackage[unicode,hidelinks,bookmarks=false]{hyperref}
\newcommand{\ie}{i.e.\ }
\newcommand{\cf}{cf.\ }
\newcommand{\resp}{respectively\ }
\newcommand{\Subsection}{\subsection}
\providecommand{\nobreakdash}{\nobreak\hbox{-}}
\setlength{\emergencystretch}{2em}
\multlinegap0pt
\usepackage{amssymb}
\usepackage[shortlabels]{enumitem}
\usepackage{natbib}
\newtheorem{lemma}{Lemma}
\title{The inverse eigenvalue problems for perturbed Bessel operator with mixed data}
\author{Zeguang Liu, Xin-Jian Xu}
\date{Source: arXiv:2601.01093v1 (CC BY 4.0)}
\begin{document}
\maketitle

\noindent\emph{Source and licence.} The inverse eigenvalue problems for perturbed Bessel operator with mixed data. Version arXiv:2601.01093v1 (CC BY 4.0), distributed by arXiv under the Creative Commons Attribution 4.0 International licence, http://creativecommons.org/licenses/by/4.0/, as recorded in the arXiv licence metadata for this exact version and shown on the abstract page https://arxiv.org/abs/2601.01093v1. Full licence notice: \texttt{licenses/arxiv-2601.01093-CC-BY-4.0.txt}.\par\smallskip

\noindent\textit{Editorial scope.} Counted unit: the article's lemma l3.6 (source0.tex lines 509--528). The article's complete original proof of it is delivered with it (source0.tex lines 530--611, one \texttt{proof} environment of 82 lines). No mathematics is written, repaired or re-derived: the statement and the proof are the article's own lines, in the article's own notation, and no retained line is substituted or re-flowed.  Retained uncounted as the source-local context the proof needs: lines 353--360; lines 419--422.  Nothing was omitted from any retained span, so the package carries no omission record.  No retained line contains a citation, so the wrapper carries no bibliography and no bibliography entry.  No retained line contains a figure environment, an image-inclusion command or drawing code, so no image is delivered and the package declares no asset dependency.  Editorial formatting only, in addition: an AMS \texttt{amsart} preamble that restates the article's own declarations and macros used by the retained lines, namely the environments \texttt{lemma} on separate counters, together with the standard packages those lines need.  The retained environments print in this document as Lemma 1 in order of appearance, whereas the article prints the same items with the numbering of its own full document.  The complete native source of the article as distributed in the arXiv e-print for this exact version is bundled unchanged as \texttt{sources/192049/source0.tex}.
\begin{align}
\phi_{\ell}(\lambda, x, q)
&=\lambda^{-\frac{\ell+1}{2}}
\left(\sin \left(\sqrt{\lambda} x-\frac{\ell \pi}{2}\right)+O\left((\textnormal{R}(\lambda)+|\lambda|^{-\frac{1}{2}}) e^{x| \textnormal{Im}(\sqrt{\lambda}) \mid}\right)\right),\label{e1.5}\\
\phi_{\ell}^{\prime}(\lambda, x, q)
&=\lambda^{-\frac{\ell}{2}}
\left(\cos \left(\sqrt{\lambda} x-\frac{\ell \pi}{2}\right)+O\left((\textnormal{R}(\lambda)+|\lambda|^{-\frac{1}{2}}) e^{x|\textnormal{Im}(\sqrt{\lambda})|}\right)\right),\label{e1.6}
\end{align}

\begin{align}\label{e3.1}
H_{\ell}(\lambda)
=\int_0^a(q(x)-\hat{q}(x))\phi_{\ell}(\lambda, x, q)\phi_{\ell}(\lambda, x, \hat{q})dx. 
\end{align}

\begin{lemma}\label{l3.6}
Suppose $\ell\geq-1/2$, $q$, $\hat{q}\in L^{1}(0, a)$ satisfying $q=\hat{q}$ a.e. on $(a, 1)$, for some $0<a\leq1$. Assume further: 

For $1\leq p<\infty$ and $k\in\mathbb{N}\cup\{0\}$, there exists $\delta_0>0$ such that $q-\hat{q}\in W^{k,p}((a-\delta_0, a])$, and
$(q-\hat{q})^{(i)}(a)=0$, for $i=0,1, \cdots, k-1$.
Then for any $\varepsilon>0$, there exists $\delta(\varepsilon)>0$ such that
\begin{equation}\label{e3.7}
|H_{\ell}(z^2)| 
\leq \frac{e^{2a|\textnormal{Im}z|}}
{|z|^{2\ell+2}|\textnormal{Im}z|^{k+\frac{1}{p^{\prime}}}}
\left(\varepsilon
+C|z|^{\gamma}e^{-\delta(\varepsilon)|\textnormal{Im}z|}
+C|\textnormal{Im}z|^{k+\frac{1}{p^{\prime}}}
e^{-2\delta(\varepsilon)|\textnormal{Im}z|}\right), 
\end{equation}
where $\textnormal{Im}z \neq 0$,
and the constants $C>0$ and $0<\gamma<1$ are independent of $\varepsilon, \delta$ and $z$.

For $p=\infty$, if $q-\hat{q} \in C^{k}((a-\delta_0, a])$ and $(q-\hat{q})^{(i)}(a)=0$ for $i=0, \cdots, k$, then \eqref{e3.7} holds.
\end{lemma}

\begin{proof}
Fix $x_0\in(0, a-\delta_0)$.
By the asymptotic estimates \eqref{e1.5} and \eqref{e1.6}, we can choose a constant $0<\gamma<1$ such that 
\begin{equation}\label{l3.11}
\phi_{\ell}(z^2, x, q)
=z^{-\ell-1}\left(\sin \left(z x-\frac{\ell\pi}{2}\right)
+O(|z|^{\gamma-1}e^{x|\textnormal{Im}z|})\right),
\end{equation}
and
\begin{equation}\label{l3.12}
\phi_{\ell}^{\prime}(z^2, x, q)
=z^{-\ell}
\left(\cos \left(zx-\frac{\ell \pi}{2}\right)+O\left(|z|^{\gamma-1} e^{x|\textnormal{Im}(z)|}\right)\right),
\end{equation}
uniformly in $x\in[x_0,1]$. The identity \eqref{e3.1} can be rewritten as
\begin{equation}\label{l3.13}
\begin{split}
H_{\ell}(z^2)=&\left(\int_{0}^{x_0}+\int_{x_0}^{a}\right)(q-\hat{q})(x)\phi_{\ell}(\lambda, x, q) \phi_{\ell}(\lambda, x, \hat{q})dx\\
=&\phi_{\ell}(z^2,x_0,q)\phi_{\ell}^{\prime}(z^2,x_0,\hat{q})-\phi_{\ell}(z^2,x_0,\hat{q})\phi_{\ell}^{\prime}(z^2,x_0,q)\\
&+\int_{x_0}^{a}(q-\hat{q})(x)\phi_{\ell}(\lambda, x, q) \phi_{\ell}(\lambda, x, \hat{q})dx.
\end{split}
\end{equation}
Putting the estimates \eqref{l3.11} and \eqref{l3.12} into the identity \eqref{l3.13}, for $\delta\in(0,\delta_{0})$, one has
\begin{align}\label{e3.3}
\begin{split}
H_{\ell}(z^2)= &z^{-2 \ell-2}\int_{x_0}^{a}(q-\hat{q})(x)f_0(x)dx+z^{-2\ell-2}g(z)\\
= &z^{-2 \ell-2}\left(\int_{x_0}^{a-\delta}+\int_{a-\delta}^{a}\right)
(q-\hat{q})(x)f_0(x)dx+z^{-2\ell-2}g(z),
\end{split}
\end{align}
where $f_0(x)=O\left(e^{2x|\textnormal{Im}z|}\right)$ uniformly in $x\in[x_0,a]$, and $g(z)=O\left(|z|^{\gamma}e^{2x_0|\textnormal{Im}z|}\right)$, 
for $|z|\rightarrow\infty$. 

On the one hand, 
since $q-\hat{q}\in W^{k, p}(a-\delta_0, a]$ and $(q-\hat{q})^{(i)}(a)=0$ for $i=0, \cdots, k-1$, and $1\leq p<\infty$, 
by performing $k$ iterations of integration by parts, we have 
\begin{equation*}
\begin{split}
\int_{a-\delta}^a(q-\hat{q})(x)f_{0}(x)dx=&(-1)\int_{a-\delta}^a(q-\hat{q})^{\prime}(x)f_{1}(x)dx=\cdots\\
=&(-1)^{k}\int_{a-\delta}^a(q-\hat{q})^{(k)}(x)f_{k}(x)dx,
\end{split}
\end{equation*}
where 
\begin{align*}
f_{i+1}(x)=\int_{a-\delta}^{x} f_i(t)dt.
\end{align*}
We see by induction on $i$ that
\begin{equation}\label{l3.14}
f_i(x)
=O\left(\frac{e^{2|\textnormal{Im}z|x}}
{|\textnormal{Im}z|^{i}}\right)
\end{equation}
uniformly in $x\in[a-\delta,a]$. Note that H\"{o}lder's inequality gives 
\begin{equation}\label{l3.15}
\int_{a-\delta}^{a}
\left|(q-\hat{q})^{(k)}(x)\right|e^{2|\textnormal{Im}z|x}dx \leq C_1\|(q-\hat{q})^{(k)}\|_{L^{p}(a-\delta, a)}
\frac{e^{2|\textnormal{Im}z|a}}{|\textnormal{Im}z|^{1/p^{\prime}}},
\end{equation}
for some constant $C_1>0$. For small $\delta$, the $L^{p}$-norm of $(q-\hat{q})^{(k)}$ is small. 
The facts \eqref{l3.14} and \eqref{l3.15} show that for small $\delta$, we have 
\begin{equation}\label{e3.4}
\left|\int_{a-\delta}^{a}(q-\hat{q})(x)f_{0}(x)dx\right| 
\leq \varepsilon \frac{e^{2|\textnormal{Im}z|a}}{|\textnormal{Im}z|^{k+1/p^{\prime}}}.
\end{equation}
From the additional information, one sees that the inequality \eqref{e3.4} also holds for $p=\infty$.


On the other hand, 
\begin{equation}\label{e3.5}
\left|\int_{x_0}^{a-\delta}(q-\hat{q})(x)f_{0}(x)dx\right| 
\leq C_{2}e^{2|\textnormal{Im}z|(a-\delta)},
\end{equation}
for some constant $C_2>0$. Moreover, the function $g(z)$ in \eqref{e3.3} has the following estimate:
\begin{equation}\label{e3.6}
 |g(z)|
\leq C_{3}
\frac{e^{2|\textnormal{Im}z| a}}{|\textnormal{Im}z|^{k+1/p^{\prime}}} |z|^{\gamma}e^{-(a-x_0)|\textnormal{Im}z|}
\leq C_{3} 
\frac{e^{2|\textnormal{Im}z|a}}{|\textnormal{Im}z|^{k+1/p^{\prime}}} |z|^{\gamma}e^{-\delta|\textnormal{Im}z|},
\end{equation}
for some constant $C_3>0$. By using \eqref{e3.3}, \eqref{e3.4}, \eqref{e3.5} and \eqref{e3.6}, we complete the proof.
\end{proof}

\end{document}
